% 6.5 230 个空间群的双值表示
\section{230 个空间群的双值表示}

本节专门列出三维空间群的双值表示并按定义 1.3.7 对其分类（所用检验见 4.6 节）。它与 5.2 节类似，后者给出了单值表示的同类表格。双值表示的列表及其标记以与单值表示相同的方式在表 6.13 与 6.14 中给出。表 6.13 类比于表 5.7，表 6.14 类比于表 5.8：表 6.13 把双值表示辨认为表 5.1 的某个抽象群的表示，表 6.14 给出空间群各双值表示所用的标记。有了这两张表的注，加上 5.4 节考虑的表 5.7 与 5.8 的使用示例以及 6.4 节已考虑的例子，读者应当不难得到 230 个空间群中任何一个的基本域内任意波矢 $\mathbf{k}$ 处 $\mathbf{G}^{\mathbf{k}}$ 的双值小表示。表 6.13 中还以编码形式再次给出关于每个空间群的其他作者工作的出处。空间群表示同时用 Mulliken (1933) 记号的推广与若干作者所用 $\Gamma$ 记号的推广来标记。若计及时间反演对称可能出现的额外简并，也一并在表中指明。

以下为表 6.13（230 个空间群的双值表示，原书第 467--560 页）与表 6.14（双值表示的标记）的扫描页。

\begin{figure}[H]
\centering
\includegraphics[width=0.86\textwidth]{c6_65_p480.jpg}
\caption*{原书第 467 页}
\end{figure}
\begin{figure}[H]
\centering
\includegraphics[width=0.86\textwidth]{c6_65_p481.jpg}
\caption*{原书第 468 页}
\end{figure}
\begin{figure}[H]
\centering
\includegraphics[width=0.86\textwidth]{c6_65_p482.jpg}
\caption*{原书第 469 页}
\end{figure}
\begin{figure}[H]
\centering
\includegraphics[width=0.86\textwidth]{c6_65_p483.jpg}
\caption*{原书第 470 页}
\end{figure}
\begin{figure}[H]
\centering
\includegraphics[width=0.86\textwidth]{c6_65_p484.jpg}
\caption*{原书第 471 页}
\end{figure}
\begin{figure}[H]
\centering
\includegraphics[width=0.86\textwidth]{c6_65_p485.jpg}
\caption*{原书第 472 页}
\end{figure}
\begin{figure}[H]
\centering
\includegraphics[width=0.86\textwidth]{c6_65_p486.jpg}
\caption*{原书第 473 页}
\end{figure}
\begin{figure}[H]
\centering
\includegraphics[width=0.86\textwidth]{c6_65_p487.jpg}
\caption*{原书第 474 页}
\end{figure}
\begin{figure}[H]
\centering
\includegraphics[width=0.86\textwidth]{c6_65_p488.jpg}
\caption*{原书第 475 页}
\end{figure}
\begin{figure}[H]
\centering
\includegraphics[width=0.86\textwidth]{c6_65_p489.jpg}
\caption*{原书第 476 页}
\end{figure}
\begin{figure}[H]
\centering
\includegraphics[width=0.86\textwidth]{c6_65_p490.jpg}
\caption*{原书第 477 页}
\end{figure}
\begin{figure}[H]
\centering
\includegraphics[width=0.86\textwidth]{c6_65_p491.jpg}
\caption*{原书第 478 页}
\end{figure}
\begin{figure}[H]
\centering
\includegraphics[width=0.86\textwidth]{c6_65_p492.jpg}
\caption*{原书第 479 页}
\end{figure}
\begin{figure}[H]
\centering
\includegraphics[width=0.86\textwidth]{c6_65_p493.jpg}
\caption*{原书第 480 页}
\end{figure}
\begin{figure}[H]
\centering
\includegraphics[width=0.86\textwidth]{c6_65_p494.jpg}
\caption*{原书第 481 页}
\end{figure}
\begin{figure}[H]
\centering
\includegraphics[width=0.86\textwidth]{c6_65_p495.jpg}
\caption*{原书第 482 页}
\end{figure}
\begin{figure}[H]
\centering
\includegraphics[width=0.86\textwidth]{c6_65_p496.jpg}
\caption*{原书第 483 页}
\end{figure}
\begin{figure}[H]
\centering
\includegraphics[width=0.86\textwidth]{c6_65_p497.jpg}
\caption*{原书第 484 页}
\end{figure}
\begin{figure}[H]
\centering
\includegraphics[width=0.86\textwidth]{c6_65_p498.jpg}
\caption*{原书第 485 页}
\end{figure}
\begin{figure}[H]
\centering
\includegraphics[width=0.86\textwidth]{c6_65_p499.jpg}
\caption*{原书第 486 页}
\end{figure}
\begin{figure}[H]
\centering
\includegraphics[width=0.86\textwidth]{c6_65_p500.jpg}
\caption*{原书第 487 页}
\end{figure}
\begin{figure}[H]
\centering
\includegraphics[width=0.86\textwidth]{c6_65_p501.jpg}
\caption*{原书第 488 页}
\end{figure}
\begin{figure}[H]
\centering
\includegraphics[width=0.86\textwidth]{c6_65_p502.jpg}
\caption*{原书第 489 页}
\end{figure}
\begin{figure}[H]
\centering
\includegraphics[width=0.86\textwidth]{c6_65_p503.jpg}
\caption*{原书第 490 页}
\end{figure}
\begin{figure}[H]
\centering
\includegraphics[width=0.86\textwidth]{c6_65_p504.jpg}
\caption*{原书第 491 页}
\end{figure}
\begin{figure}[H]
\centering
\includegraphics[width=0.86\textwidth]{c6_65_p505.jpg}
\caption*{原书第 492 页}
\end{figure}
\begin{figure}[H]
\centering
\includegraphics[width=0.86\textwidth]{c6_65_p506.jpg}
\caption*{原书第 493 页}
\end{figure}
\begin{figure}[H]
\centering
\includegraphics[width=0.86\textwidth]{c6_65_p507.jpg}
\caption*{原书第 494 页}
\end{figure}
\begin{figure}[H]
\centering
\includegraphics[width=0.86\textwidth]{c6_65_p508.jpg}
\caption*{原书第 495 页}
\end{figure}
\begin{figure}[H]
\centering
\includegraphics[width=0.86\textwidth]{c6_65_p509.jpg}
\caption*{原书第 496 页}
\end{figure}
\begin{figure}[H]
\centering
\includegraphics[width=0.86\textwidth]{c6_65_p510.jpg}
\caption*{原书第 497 页}
\end{figure}
\begin{figure}[H]
\centering
\includegraphics[width=0.86\textwidth]{c6_65_p511.jpg}
\caption*{原书第 498 页}
\end{figure}
\begin{figure}[H]
\centering
\includegraphics[width=0.86\textwidth]{c6_65_p512.jpg}
\caption*{原书第 499 页}
\end{figure}
\begin{figure}[H]
\centering
\includegraphics[width=0.86\textwidth]{c6_65_p513.jpg}
\caption*{原书第 500 页}
\end{figure}
\begin{figure}[H]
\centering
\includegraphics[width=0.86\textwidth]{c6_65_p514.jpg}
\caption*{原书第 501 页}
\end{figure}
\begin{figure}[H]
\centering
\includegraphics[width=0.86\textwidth]{c6_65_p515.jpg}
\caption*{原书第 502 页}
\end{figure}
\begin{figure}[H]
\centering
\includegraphics[width=0.86\textwidth]{c6_65_p516.jpg}
\caption*{原书第 503 页}
\end{figure}
\begin{figure}[H]
\centering
\includegraphics[width=0.86\textwidth]{c6_65_p517.jpg}
\caption*{原书第 504 页}
\end{figure}
\begin{figure}[H]
\centering
\includegraphics[width=0.86\textwidth]{c6_65_p518.jpg}
\caption*{原书第 505 页}
\end{figure}
\begin{figure}[H]
\centering
\includegraphics[width=0.86\textwidth]{c6_65_p519.jpg}
\caption*{原书第 506 页}
\end{figure}
\begin{figure}[H]
\centering
\includegraphics[width=0.86\textwidth]{c6_65_p520.jpg}
\caption*{原书第 507 页}
\end{figure}
\begin{figure}[H]
\centering
\includegraphics[width=0.86\textwidth]{c6_65_p521.jpg}
\caption*{原书第 508 页}
\end{figure}
\begin{figure}[H]
\centering
\includegraphics[width=0.86\textwidth]{c6_65_p522.jpg}
\caption*{原书第 509 页}
\end{figure}
\begin{figure}[H]
\centering
\includegraphics[width=0.86\textwidth]{c6_65_p523.jpg}
\caption*{原书第 510 页}
\end{figure}
\begin{figure}[H]
\centering
\includegraphics[width=0.86\textwidth]{c6_65_p524.jpg}
\caption*{原书第 511 页}
\end{figure}
\begin{figure}[H]
\centering
\includegraphics[width=0.86\textwidth]{c6_65_p525.jpg}
\caption*{原书第 512 页}
\end{figure}
\begin{figure}[H]
\centering
\includegraphics[width=0.86\textwidth]{c6_65_p526.jpg}
\caption*{原书第 513 页}
\end{figure}
\begin{figure}[H]
\centering
\includegraphics[width=0.86\textwidth]{c6_65_p527.jpg}
\caption*{原书第 514 页}
\end{figure}
\begin{figure}[H]
\centering
\includegraphics[width=0.86\textwidth]{c6_65_p528.jpg}
\caption*{原书第 515 页}
\end{figure}
\begin{figure}[H]
\centering
\includegraphics[width=0.86\textwidth]{c6_65_p529.jpg}
\caption*{原书第 516 页}
\end{figure}
\begin{figure}[H]
\centering
\includegraphics[width=0.86\textwidth]{c6_65_p530.jpg}
\caption*{原书第 517 页}
\end{figure}
\begin{figure}[H]
\centering
\includegraphics[width=0.86\textwidth]{c6_65_p531.jpg}
\caption*{原书第 518 页}
\end{figure}
\begin{figure}[H]
\centering
\includegraphics[width=0.86\textwidth]{c6_65_p532.jpg}
\caption*{原书第 519 页}
\end{figure}
\begin{figure}[H]
\centering
\includegraphics[width=0.86\textwidth]{c6_65_p533.jpg}
\caption*{原书第 520 页}
\end{figure}
\begin{figure}[H]
\centering
\includegraphics[width=0.86\textwidth]{c6_65_p534.jpg}
\caption*{原书第 521 页}
\end{figure}
\begin{figure}[H]
\centering
\includegraphics[width=0.86\textwidth]{c6_65_p535.jpg}
\caption*{原书第 522 页}
\end{figure}
\begin{figure}[H]
\centering
\includegraphics[width=0.86\textwidth]{c6_65_p536.jpg}
\caption*{原书第 523 页}
\end{figure}
\begin{figure}[H]
\centering
\includegraphics[width=0.86\textwidth]{c6_65_p537.jpg}
\caption*{原书第 524 页}
\end{figure}
\begin{figure}[H]
\centering
\includegraphics[width=0.86\textwidth]{c6_65_p538.jpg}
\caption*{原书第 525 页}
\end{figure}
\begin{figure}[H]
\centering
\includegraphics[width=0.86\textwidth]{c6_65_p539.jpg}
\caption*{原书第 526 页}
\end{figure}
\begin{figure}[H]
\centering
\includegraphics[width=0.86\textwidth]{c6_65_p540.jpg}
\caption*{原书第 527 页}
\end{figure}
\begin{figure}[H]
\centering
\includegraphics[width=0.86\textwidth]{c6_65_p541.jpg}
\caption*{原书第 528 页}
\end{figure}
\begin{figure}[H]
\centering
\includegraphics[width=0.86\textwidth]{c6_65_p542.jpg}
\caption*{原书第 529 页}
\end{figure}
\begin{figure}[H]
\centering
\includegraphics[width=0.86\textwidth]{c6_65_p543.jpg}
\caption*{原书第 530 页}
\end{figure}
\begin{figure}[H]
\centering
\includegraphics[width=0.86\textwidth]{c6_65_p544.jpg}
\caption*{原书第 531 页}
\end{figure}
\begin{figure}[H]
\centering
\includegraphics[width=0.86\textwidth]{c6_65_p545.jpg}
\caption*{原书第 532 页}
\end{figure}
\begin{figure}[H]
\centering
\includegraphics[width=0.86\textwidth]{c6_65_p546.jpg}
\caption*{原书第 533 页}
\end{figure}
\begin{figure}[H]
\centering
\includegraphics[width=0.86\textwidth]{c6_65_p547.jpg}
\caption*{原书第 534 页}
\end{figure}
\begin{figure}[H]
\centering
\includegraphics[width=0.86\textwidth]{c6_65_p548.jpg}
\caption*{原书第 535 页}
\end{figure}
\begin{figure}[H]
\centering
\includegraphics[width=0.86\textwidth]{c6_65_p549.jpg}
\caption*{原书第 536 页}
\end{figure}
\begin{figure}[H]
\centering
\includegraphics[width=0.86\textwidth]{c6_65_p550.jpg}
\caption*{原书第 537 页}
\end{figure}
\begin{figure}[H]
\centering
\includegraphics[width=0.86\textwidth]{c6_65_p551.jpg}
\caption*{原书第 538 页}
\end{figure}
\begin{figure}[H]
\centering
\includegraphics[width=0.86\textwidth]{c6_65_p552.jpg}
\caption*{原书第 539 页}
\end{figure}
\begin{figure}[H]
\centering
\includegraphics[width=0.86\textwidth]{c6_65_p553.jpg}
\caption*{原书第 540 页}
\end{figure}
\begin{figure}[H]
\centering
\includegraphics[width=0.86\textwidth]{c6_65_p554.jpg}
\caption*{原书第 541 页}
\end{figure}
\begin{figure}[H]
\centering
\includegraphics[width=0.86\textwidth]{c6_65_p555.jpg}
\caption*{原书第 542 页}
\end{figure}
\begin{figure}[H]
\centering
\includegraphics[width=0.86\textwidth]{c6_65_p556.jpg}
\caption*{原书第 543 页}
\end{figure}
\begin{figure}[H]
\centering
\includegraphics[width=0.86\textwidth]{c6_65_p557.jpg}
\caption*{原书第 544 页}
\end{figure}
\begin{figure}[H]
\centering
\includegraphics[width=0.86\textwidth]{c6_65_p558.jpg}
\caption*{原书第 545 页}
\end{figure}
\begin{figure}[H]
\centering
\includegraphics[width=0.86\textwidth]{c6_65_p559.jpg}
\caption*{原书第 546 页}
\end{figure}
\begin{figure}[H]
\centering
\includegraphics[width=0.86\textwidth]{c6_65_p560.jpg}
\caption*{原书第 547 页}
\end{figure}
\begin{figure}[H]
\centering
\includegraphics[width=0.86\textwidth]{c6_65_p561.jpg}
\caption*{原书第 548 页}
\end{figure}
\begin{figure}[H]
\centering
\includegraphics[width=0.86\textwidth]{c6_65_p562.jpg}
\caption*{原书第 549 页}
\end{figure}
\begin{figure}[H]
\centering
\includegraphics[width=0.86\textwidth]{c6_65_p563.jpg}
\caption*{原书第 550 页}
\end{figure}
\begin{figure}[H]
\centering
\includegraphics[width=0.86\textwidth]{c6_65_p564.jpg}
\caption*{原书第 551 页}
\end{figure}
\begin{figure}[H]
\centering
\includegraphics[width=0.86\textwidth]{c6_65_p565.jpg}
\caption*{原书第 552 页}
\end{figure}
\begin{figure}[H]
\centering
\includegraphics[width=0.86\textwidth]{c6_65_p566.jpg}
\caption*{原书第 553 页}
\end{figure}
\begin{figure}[H]
\centering
\includegraphics[width=0.86\textwidth]{c6_65_p567.jpg}
\caption*{原书第 554 页}
\end{figure}
\begin{figure}[H]
\centering
\includegraphics[width=0.86\textwidth]{c6_65_p568.jpg}
\caption*{原书第 555 页}
\end{figure}
\begin{figure}[H]
\centering
\includegraphics[width=0.86\textwidth]{c6_65_p569.jpg}
\caption*{原书第 556 页}
\end{figure}
\begin{figure}[H]
\centering
\includegraphics[width=0.86\textwidth]{c6_65_p570.jpg}
\caption*{原书第 557 页}
\end{figure}
\begin{figure}[H]
\centering
\includegraphics[width=0.86\textwidth]{c6_65_p571.jpg}
\caption*{原书第 558 页}
\end{figure}
\begin{figure}[H]
\centering
\includegraphics[width=0.86\textwidth]{c6_65_p572.jpg}
\caption*{原书第 559 页}
\end{figure}
\begin{figure}[H]
\centering
\includegraphics[width=0.86\textwidth]{c6_65_p573.jpg}
\caption*{原书第 560 页}
\end{figure}
